% 题干转写：完成
% 解析转写：完成
% 图形核对：完成
\section{第37届复赛第3题：负微分电阻RLC振荡电路}
\subsection{题干}
隧道二极管在特定电压范围内具有负微分电阻 $-R_0$（$R_0>0$）。如图\ref{fig:cpho37q03}，左侧回路由定值电阻 $R$、电感 $L$、电容 $C$ 和二极管 $D$ 组成；右侧由理想高频扼流圈与理想恒压直流源 $V_0$ 组成。设左侧交流电流为 $i(t)$，且 $i(0)=0$、$\dot i(0)=\beta$。求：（1）$i(t)$ 的一般表达式；（2）谐振条件、谐振电流与频率；（3）阻尼振荡条件、电流与频率；（4）在工作电压范围 $V_{\min}\le V_D\le V_{\max}$ 内，二极管交流部分的最大平均功率及此时 $V_0$。
\begin{figure}[htbp]\centering
\begin{tikzpicture}[scale=0.95,>=Stealth,font=\small]
  % AC RLC loop with a negative differential-resistance tunnel diode.
  \draw[thick] (0,0) -- (0,2.0);
  \draw[thick] (0,2.0) -- (0.55,2.0);
  \draw[thick,decorate,decoration={coil,aspect=0.55,segment length=4.5pt,amplitude=3.8pt}] (0.55,2.0) -- (1.85,2.0);
  \node[above] at (1.18,2.18) {$L$};
  \draw[thick] (1.85,2.0) -- (2.18,2.0);
  \draw[thick] (2.25,1.66) -- (2.25,2.34);
  \draw[thick] (2.48,1.66) -- (2.48,2.34);
  \draw[thick] (2.18,2.0) -- (2.25,2.0) (2.48,2.0) -- (3.0,2.0);
  \node[above] at (2.36,2.30) {$C$};
  \draw[thick] (3.0,2.0) -- (3.0,1.20);
  \draw[fill=gray!15,thick] (2.72,0.65) rectangle (3.28,1.20);
  \node at (3.0,0.92) {$D$};
  \draw[thick] (3.0,0.65) -- (3.0,0) -- (0,0);
  \draw[thick] (0,1.54) -- (-0.20,1.54);
  \draw[thick,decorate,decoration={zigzag,segment length=4pt,amplitude=1.8pt}] (-0.20,1.54) -- (-0.20,0.46);
  \draw[thick] (-0.20,0.46) -- (0,0.46);
  \node[left] at (-0.20,1.02) {$R$};
  \draw[->,thick] (0.22,1.05) arc[start angle=155,end angle=28,radius=0.95] node[midway,above left] {$i(t)$};
  \draw[<->] (3.48,0.02) -- (3.48,1.18) node[midway,right] {$V_D$};

  % DC branch isolated from the AC loop by an RF choke.
  \draw[thick] (3.0,2.0) -- (4.0,2.0);
  \draw[thick,decorate,decoration={coil,aspect=0.55,segment length=3.8pt,amplitude=3.0pt}] (4.0,2.0) -- (5.10,2.0);
  \node[above] at (4.55,2.18) {RFC};
  \draw[thick] (5.10,2.0) -- (5.60,2.0) -- (5.60,0) -- (3.0,0);
  \draw[thick] (5.42,0.95) -- (5.78,0.95);
  \draw[thick] (5.49,0.63) -- (5.71,0.63);
  \node[right] at (5.78,0.77) {$V_0$};
  \draw[->] (4.30,1.70) -- (5.05,1.70) node[midway,below] {$I$};
\end{tikzpicture}

\caption{第37届复赛第3题负微分电阻RLC电路的 TikZ 重绘。}\label{fig:cpho37q03}
\end{figure}

\subsection{参考解答整理}
\paragraph{（1）一般电流表达式。} 左侧回路满足
\[
Ri(t)+L\dot i(t)+\frac1C\int_0^t i(t')\,\mathrm dt'-R_0i(t)=0.
\]
微分后得到特征方程
\[
L\lambda^2+(R-R_0)\lambda+\frac1C=0.
\]
其根为
\[
\lambda_\pm=-\frac{R-R_0}{2L}
\pm\sqrt{\left(\frac{R-R_0}{2L}\right)^2-\frac1{LC}}.
\]
由初始条件，电流写为
\[
\boxed{i(t)=\frac{\beta}{2\sqrt{\left[(R-R_0)/(2L)\right]^2-1/(LC)}}
\left(e^{\lambda_+t}-e^{\lambda_-t}\right).}
\]

\paragraph{（2）谐振。} 要使能量不随时间耗损且电流周期变化，应有
\[
R=R_0.
\]
此时
\[
\boxed{i_H(t)=\beta\sqrt{LC}\sin\frac{t}{\sqrt{LC}},
\qquad f_H=\frac1{2\pi\sqrt{LC}}.}
\]

\paragraph{（3）RLC阻尼振荡。} 能量衰减且仍保持周期变化的条件为
\[
R-R_0>0,
\qquad
\left(\frac{R-R_0}{2L}\right)^2<\frac1{LC},
\]
即
\[
\boxed{R_0<R<R_0+2\sqrt{\frac LC}.}
\]
电流及频率分别为
\[
\boxed{i_D(t)=\frac{\beta}{\sqrt{1/(LC)-[(R-R_0)/(2L)]^2}}
e^{-(R-R_0)t/(2L)}
\sin\left(\sqrt{\frac1{LC}-\left[\frac{R-R_0}{2L}\right]^2}\,t\right),}
\]
\[
\boxed{f_D=\frac1{2\pi}\sqrt{\frac1{LC}-\left[\frac{R-R_0}{2L}\right]^2}.}
\]

\paragraph{（4）最大平均功率。} 为在负微分电阻区内获得最大交流电压幅值，直流偏置应取工作范围中点：
\[
\boxed{V_0=\bar V=\frac{V_{\max}+V_{\min}}2,\qquad
V_{D\max}=\frac{V_{\max}-V_{\min}}2.}
\]
因此二极管交流部分的最大平均功率为
\[
\boxed{\bar P_{\max}=\frac12\frac{V_{D\max}^2}{R_0}
=\frac{(V_{\max}-V_{\min})^2}{8R_0}.}
\]

\subsection{转写依据}
题干、答案公式和图形分别依据 \texttt{CPHO37-Q03.md}、\texttt{CPHO37-Q03\_formula\_check\_v1.csv} 与题图核对入口。评分分布为第（1）问18分、第（2）问7分、第（3）问9分、第（4）问6分。
